\chapter{Appendicitis Tests}

\section*{What Is This Test?}
Appendicitis testing is a clinical workup used to determine whether the appendix is inflamed and whether complications such as perforation or an abscess are present. No single blood or urine test proves or excludes appendicitis.

\section*{Alternative Names}
\begin{itemize}
  \item Appendicitis Workup
  \item Appendix Evaluation
  \item Suspected Acute Appendicitis Testing
\end{itemize}

\section*{Principle}
The diagnosis combines the history and abdominal examination with risk scores, laboratory tests, and imaging. A complete blood count and inflammatory markers may support the diagnosis, while ultrasound, CT, or MRI can show an enlarged or inflamed appendix and complications.

\section*{Why This Test Is Ordered}
It is ordered for right-lower-quadrant pain, migrating abdominal pain, nausea or vomiting, fever, guarding, or other symptoms concerning for acute appendicitis. Pregnancy, age, body habitus, and alternative diagnoses affect the choice of imaging.

\section*{Primary vs Secondary Test}
Clinical assessment and appropriate imaging are primary components. CBC, CRP, urinalysis, electrolytes, and a pregnancy test are supportive or safety tests. A normal white-cell count does not reliably exclude appendicitis.

\section*{How This Test Is Performed}
Blood and urine are collected when indicated. Ultrasound is often preferred initially in children and pregnancy; CT provides detailed assessment in many adults; MRI is an option when avoiding ionizing radiation is important and it is available.

\section*{What Preparation Is Needed}
Do not eat or drink until the treating team advises otherwise if surgery or a procedure may be needed. Tell the clinician about pregnancy possibility, medicines, allergies, and prior abdominal surgery.

\section*{Contraindications and Precautions}
Imaging choice should account for radiation exposure, contrast allergy, kidney function, and pregnancy. Severe or worsening pain, rigid abdomen, fainting, persistent vomiting, or sepsis requires urgent evaluation.

\section*{Understanding Results}
Leukocytosis or raised CRP supports inflammation but is nonspecific. Imaging findings must be interpreted with symptoms and examination. An abscess, perforation, obstruction, or diffuse peritonitis changes urgency and management.

\section*{Follow-up Tests}
Depending on the findings, follow-up may include repeat examination and blood tests, CT or ultrasound, surgical consultation, pathology examination of a removed appendix, and cultures when perforation or abscess is present.

\section*{References}
\begin{enumerate}
  \item MedlinePlus. Appendicitis Tests.
  \item American College of Radiology. Appropriateness Criteria: Right Lower Quadrant Pain.
  \item World Society of Emergency Surgery. Jerusalem guidelines for acute appendicitis.
\end{enumerate}

\clearpage
\section*{Understanding Results and Follow-up}
The laboratory report should identify the specimen, method, units, reference interval or decision limit, and whether the result is positive, negative, abnormal, or inconclusive. Reference intervals vary with age, sex, pregnancy, specimen type, assay, and laboratory; a result outside a range is not automatically a diagnosis, and a result within a range does not always exclude disease. Discuss unexpected results with the ordering clinician, who can decide whether repeat or confirmatory testing, treatment, or referral is appropriate. Do not change medicines or delay urgent care based on an isolated result.


\section*{Regulatory Status and Authorization Date}
This chapter describes a test category, not a specific commercial product. FDA, CDSCO, EU IVDR, and MHRA approval or registration dates therefore cannot be assigned without the assay manufacturer, product name, instrument, and intended use. For laboratory-developed tests, an individual agency approval date may not exist. See the Preface for the interpretation of regulatory status.

\clearpage
